% Options for packages loaded elsewhere
\PassOptionsToPackage{unicode}{hyperref}
\PassOptionsToPackage{hyphens}{url}
\documentclass[
]{article}
\usepackage{xcolor}
\usepackage[margin=1in]{geometry}
\usepackage{amsmath,amssymb}
\setcounter{secnumdepth}{-\maxdimen} % remove section numbering
\usepackage{iftex}
\ifPDFTeX
  \usepackage[T1]{fontenc}
  \usepackage[utf8]{inputenc}
  \usepackage{textcomp} % provide euro and other symbols
\else % if luatex or xetex
  \usepackage{unicode-math} % this also loads fontspec
  \defaultfontfeatures{Scale=MatchLowercase}
  \defaultfontfeatures[\rmfamily]{Ligatures=TeX,Scale=1}
\fi
\usepackage{lmodern}
\ifPDFTeX\else
  % xetex/luatex font selection
\fi
% Use upquote if available, for straight quotes in verbatim environments
\IfFileExists{upquote.sty}{\usepackage{upquote}}{}
\IfFileExists{microtype.sty}{% use microtype if available
  \usepackage[]{microtype}
  \UseMicrotypeSet[protrusion]{basicmath} % disable protrusion for tt fonts
}{}
\makeatletter
\@ifundefined{KOMAClassName}{% if non-KOMA class
  \IfFileExists{parskip.sty}{%
    \usepackage{parskip}
  }{% else
    \setlength{\parindent}{0pt}
    \setlength{\parskip}{6pt plus 2pt minus 1pt}}
}{% if KOMA class
  \KOMAoptions{parskip=half}}
\makeatother
\usepackage{longtable,booktabs,array}
\newcounter{none} % for unnumbered tables
\usepackage{calc} % for calculating minipage widths
% Correct order of tables after \paragraph or \subparagraph
\usepackage{etoolbox}
\makeatletter
\patchcmd\longtable{\par}{\if@noskipsec\mbox{}\fi\par}{}{}
\makeatother
% Allow footnotes in longtable head/foot
\IfFileExists{footnotehyper.sty}{\usepackage{footnotehyper}}{\usepackage{footnote}}
\makesavenoteenv{longtable}
\usepackage{graphicx}
\makeatletter
\newsavebox\pandoc@box
\newcommand*\pandocbounded[1]{% scales image to fit in text height/width
  \sbox\pandoc@box{#1}%
  \Gscale@div\@tempa{\textheight}{\dimexpr\ht\pandoc@box+\dp\pandoc@box\relax}%
  \Gscale@div\@tempb{\linewidth}{\wd\pandoc@box}%
  \ifdim\@tempb\p@<\@tempa\p@\let\@tempa\@tempb\fi% select the smaller of both
  \ifdim\@tempa\p@<\p@\scalebox{\@tempa}{\usebox\pandoc@box}%
  \else\usebox{\pandoc@box}%
  \fi%
}
% Set default figure placement to htbp
\def\fps@figure{htbp}
\makeatother
\setlength{\emergencystretch}{3em} % prevent overfull lines
\providecommand{\tightlist}{%
  \setlength{\itemsep}{0pt}\setlength{\parskip}{0pt}}
\usepackage{bookmark}
\IfFileExists{xurl.sty}{\usepackage{xurl}}{} % add URL line breaks if available
\urlstyle{same}
\hypersetup{
  hidelinks,
  pdfcreator={LaTeX via pandoc}}

\author{}
\date{\vspace{-2.5em}}

\begin{document}

\section{Module 4 --- Variant Calling:
Runbook}\label{module-4-variant-calling-runbook}

This is an \textbf{operational} document --- step-by-step execution
instructions, expected outputs, and troubleshooting. For the conceptual
``why,'' see this module's \href{README.md}{\texttt{README.md}}. For the
full curriculum context, see
\href{../../docs/curriculum/modular_bioinformatics_curriculum.md}{\texttt{docs/curriculum/modular\_bioinformatics\_curriculum.md}}.

\subsection{Prerequisites}\label{prerequisites}

\begin{itemize}
\tightlist
\item
  Slurm environment with \texttt{module\ load} access to
  \texttt{SRA-Toolkit/3.2.0-gompi-2024a},
  \texttt{fastp/0.23.4-GCC-13.2.0}, \texttt{BWA/0.7.18-GCC-13.3.0},
  \texttt{SAMtools/1.21-GCC-13.3.0}, \texttt{BCFtools/1.21-GCC-13.3.0},
  and \texttt{MultiQC/1.28-foss-2024a} --- used by Tier A
  (\texttt{master\_snp\_pipeline.sh}).
\item
  \texttt{samtools} and \texttt{bcftools} on \texttt{PATH} (outside
  Slurm too) for Tier B Scripts 1, 2, 3, 4, and 6 --- these are plain
  shell scripts, not Slurm jobs, and are meant to be run interactively
  against Tier A's output.
\item
  R with \texttt{vcfR}, \texttt{adegenet}, and \texttt{ape} installed
  --- used by Tier A's \texttt{variant\_analysis.R} and Tier B Scripts 5
  and 7.
\item
  Module 1 completed or reviewed --- this module assumes the same
  ``validate/interpret before you trust it'' habit taught there, applied
  one level up the pipeline (aligned reads and called variants instead
  of raw reads).
\end{itemize}

\subsection{A real, verified data-provenance issue (read before running
the smoke
test)}\label{a-real-verified-data-provenance-issue-read-before-running-the-smoke-test}

\texttt{master\_snp\_pipeline.sh}'s own built-in fallback
\texttt{srr\_list.txt} (used only when no \texttt{srr\_list.txt} is
staged) is \texttt{SRR11092056}/\texttt{SRR11092057}. Verified via NCBI
eutils: \texttt{SRR11092056} is SARS-CoV-2 metagenomic RNA-seq (Wuhan
Institute of Virology), not bacterial WGS, and cannot produce a
meaningful alignment against this pipeline's Typhimurium LT2 reference.
\textbf{Do not rely on the built-in fallback.} Use
\texttt{smoke\_test\_prjna242847.sh}, which stages a real, verified
3-isolate Typhimurium cohort (PRJNA242847) before invoking
\texttt{master\_snp\_pipeline.sh}. See the README's data-provenance note
for the full verification trail.

\subsection{Execution order}\label{execution-order}

{\def\LTcaptype{none} % do not increment counter
\begin{longtable}[]{@{}
  >{\raggedright\arraybackslash}p{(\linewidth - 4\tabcolsep) * \real{0.3333}}
  >{\raggedright\arraybackslash}p{(\linewidth - 4\tabcolsep) * \real{0.3333}}
  >{\raggedright\arraybackslash}p{(\linewidth - 4\tabcolsep) * \real{0.3333}}@{}}
\toprule\noalign{}
\begin{minipage}[b]{\linewidth}\raggedright
Step
\end{minipage} & \begin{minipage}[b]{\linewidth}\raggedright
Command
\end{minipage} & \begin{minipage}[b]{\linewidth}\raggedright
Expected result
\end{minipage} \\
\midrule\noalign{}
\endhead
\bottomrule\noalign{}
\endlastfoot
1 & \texttt{bash\ smoke\_test\_prjna242847.sh} (or
\texttt{sbatch\ master\_snp\_pipeline.sh} with your own verified
\texttt{srr\_list.txt}) &
\texttt{variants/\textless{}SRR\textgreater{}\_raw.vcf},
\texttt{variants/\textless{}SRR\textgreater{}\_filtered.vcf},
\texttt{variants/merged\_cohort.vcf} \\
2 &
\texttt{bash\ 1\_Validate\_Alignment\_And\_VCF\_Inputs.sh\ {[}WORKDIR{]}}
& Console PASS/FAIL/WARN report per sample; stops hard on a structural
or name-mismatch failure \\
3 &
\texttt{bash\ 2\_Assess\_Mapping\_And\_Reference\_Bias.sh\ {[}WORKDIR{]}}
& \texttt{mapping\_bias\_summary.tsv}; console table of \% mapped, \%
properly paired, mean depth, \% zero-coverage per sample \\
4 &
\texttt{bash\ 3\_Design\_Defensible\_Variant\_Filter.sh\ {[}WORKDIR{]}}
& \texttt{variants/\textless{}SRR\textgreater{}\_filtered\_v2.vcf},
\texttt{variants/\textless{}SRR\textgreater{}\_strand\_biased\_sites.tsv};
console naive-vs-v2 PASS-count comparison \\
5 & \texttt{bash\ 4\_Check\_Ploidy\_Assumptions.sh\ {[}WORKDIR{]}} &
\texttt{variants/\textless{}SRR\textgreater{}\_raw\_haploid.vcf};
console before/after heterozygous-call count per sample \\
6 & Re-merge the ploidy-corrected calls:
\texttt{cd\ WORKDIR/variants\ \&\&\ for\ f\ in\ *\_raw\_haploid.vcf;\ do\ bcftools\ view\ -Oz\ -o\ \$\{f\%.vcf\}.vcf.gz\ \$f\ \&\&\ bcftools\ index\ -t\ \$\{f\%.vcf\}.vcf.gz;\ done\ \&\&\ bcftools\ merge\ *\_raw\_haploid.vcf.gz\ -Ov\ -o\ merged\_cohort\_haploid.vcf}
& \texttt{variants/merged\_cohort\_haploid.vcf} --- the ploidy-correct
cohort VCF \\
7 &
\texttt{Rscript\ 5\_Cohort\_VCF\_QC\_Summary.R\ WORKDIR/variants/merged\_cohort\_haploid.vcf}
& \texttt{cohort\_vcf\_qc\_summary.tsv},
\texttt{cohort\_vcf\_qc\_per\_sample.tsv}; console Ti/Tv and
heterozygosity report \\
8 & \texttt{bash\ 6\_Coverage\_Gaps\_And\_Limitations.sh\ {[}WORKDIR{]}}
& \texttt{LIMITATIONS.md} \\
9 & \texttt{cd\ WORKDIR/variants\ \&\&\ Rscript\ variant\_analysis.R}
(Tier A, against \texttt{merged\_cohort\_haploid.vcf} --- see script
header for the required \texttt{bcftools\ merge} setup if not reusing
step 6's output) & \texttt{variant\_quality\_qc.pdf},
\texttt{phylogenetic\_tree.pdf}, \texttt{outbreak\_clusters.tsv} (if any
pair is within \texttt{SNP\_THRESHOLD}) \\
10 & \texttt{Rscript\ 7\_Assemble\_Variant\_Report.R\ WORKDIR} &
\texttt{VARIANT\_REPORT.md}, sections 1-7 auto-filled, section 8
blank \\
11 & Fill in Section 8 (interpretation) of \texttt{VARIANT\_REPORT.md}
by hand & Completed portfolio artifact \\
\end{longtable}
}

Step 6 is the one manual hand-off in this module: Script 4 corrects the
\emph{calling} step but does not silently rewrite Tier A's cohort merge
for you --- running Scripts 5-7 against the original
\texttt{merged\_cohort.vcf} instead of the ploidy-corrected re-merge
will produce numbers Script 5 itself will flag (non-zero heterozygosity
against a haploid reference), not numbers that are silently wrong.

\subsection{Expected outputs, by step}\label{expected-outputs-by-step}

\begin{itemize}
\tightlist
\item
  \textbf{Step 1 (Tier A):}
  \texttt{alignment/\textless{}SRR\textgreater{}\_dedup.bam{[}.bai{]}},
  \texttt{variants/\textless{}SRR\textgreater{}\_raw.vcf},
  \texttt{variants/\textless{}SRR\textgreater{}\_filtered.vcf},
  \texttt{variants/merged\_cohort.vcf},
  \texttt{final\_multiqc\_report.html}.
\item
  \textbf{Step 2:} stdout only --- validation report.
\item
  \textbf{Step 3:} \texttt{mapping\_bias\_summary.tsv}.
\item
  \textbf{Step 4:}
  \texttt{variants/\textless{}SRR\textgreater{}\_filtered\_v2.vcf},
  \texttt{variants/\textless{}SRR\textgreater{}\_strand\_biased\_sites.tsv},
  \texttt{variants/\textless{}SRR\textgreater{}\_raw\_metrics.tsv}.
\item
  \textbf{Step 5:}
  \texttt{variants/\textless{}SRR\textgreater{}\_raw\_haploid.vcf}.
\item
  \textbf{Step 6:} \texttt{variants/merged\_cohort\_haploid.vcf}.
\item
  \textbf{Step 7:} \texttt{cohort\_vcf\_qc\_summary.tsv},
  \texttt{cohort\_vcf\_qc\_per\_sample.tsv}.
\item
  \textbf{Step 8:} \texttt{LIMITATIONS.md}.
\item
  \textbf{Step 9 (Tier A):} \texttt{variant\_quality\_qc.pdf},
  \texttt{phylogenetic\_tree.pdf}, \texttt{outbreak\_clusters.tsv}
  (conditional).
\item
  \textbf{Step 10:} \texttt{VARIANT\_REPORT.md}.
\end{itemize}

\subsection{Troubleshooting}\label{troubleshooting}

{\def\LTcaptype{none} % do not increment counter
\begin{longtable}[]{@{}
  >{\raggedright\arraybackslash}p{(\linewidth - 4\tabcolsep) * \real{0.3333}}
  >{\raggedright\arraybackslash}p{(\linewidth - 4\tabcolsep) * \real{0.3333}}
  >{\raggedright\arraybackslash}p{(\linewidth - 4\tabcolsep) * \real{0.3333}}@{}}
\toprule\noalign{}
\begin{minipage}[b]{\linewidth}\raggedright
Symptom
\end{minipage} & \begin{minipage}[b]{\linewidth}\raggedright
Likely cause
\end{minipage} & \begin{minipage}[b]{\linewidth}\raggedright
Fix
\end{minipage} \\
\midrule\noalign{}
\endhead
\bottomrule\noalign{}
\endlastfoot
\texttt{smoke\_test\_prjna242847.sh} fails at the
\texttt{prefetch}/\texttt{fasterq-dump} step & SRA-Toolkit module not
loaded, or a transient NCBI connectivity issue & Confirm
\texttt{module\ load\ SRA-Toolkit/3.2.0-gompi-2024a} succeeded; retry
--- \texttt{master\_snp\_pipeline.sh}'s per-sample loop is
fault-tolerant and will skip a genuinely failed sample rather than abort
the whole run \\
Script 1 fails with
\texttt{@RG\ SM:\ tag\ does\ not\ match\ filename-derived\ sample} & A
BAM was renamed/copied after alignment without regenerating its read
group, or two samples' outputs were cross-copied & Re-run alignment for
that sample from \texttt{master\_snp\_pipeline.sh} rather than renaming
files by hand --- the
\texttt{-R\ "@RG\textbackslash{}tID:...\textbackslash{}tSM:..."} flag is
what makes this checkable at all \\
Script 2 flags most/all samples \texttt{LOW\_MAPPING\_RATE} & You are
accidentally running this against the pipeline's own broken default
accessions (SARS-CoV-2 RNA-seq vs.~a bacterial DNA reference) instead of
a verified worked-example cohort & Confirm \texttt{srr\_list.txt}
matches \texttt{smoke\_test\_prjna242847.sh}'s accessions or your own
verified bacterial WGS list --- see the README's data-provenance note \\
Script 3 reports 0 strand-biased sites for every sample & Either
genuinely clean data, or too few candidate calls with \texttt{INFO/DP4}
ALT depth \textgreater= 4 to evaluate (very shallow coverage) &
Cross-check against Script 2's mean-depth column before treating ``0
flagged'' as ``no strand bias present'' \\
Script 4 reports 0 heterozygous calls under the diploid default & Does
not confirm ploidy was handled correctly --- it means this particular
run got lucky, not that the risk doesn't exist. The
\texttt{*\_raw\_haploid.vcf} files are still the correct calls to build
on & Always continue with the haploid-corrected files regardless of
whether this run happened to show 0 het calls \\
Script 5 reports non-zero heterozygosity & You ran it against Tier A's
original \texttt{merged\_cohort.vcf} (diploid-default calls) instead of
the ploidy-corrected re-merge from Step 6 & Re-merge
\texttt{*\_raw\_haploid.vcf} per Step 6 and re-run Script 5 against that
file \\
Script 5 reports Ti/Tv \textless= 0.6 & Cohort's SNP calls skew toward
random noise rather than real biological signal --- could be genuinely
low-quality data, or evidence you are still running against the wrong
(viral RNA-seq vs.~bacterial reference) dataset & Confirm the dataset
first (see the mapping-rate check above); if the dataset is correct,
treat this as a real quality warning, not something to filter past
silently \\
\texttt{variant\_analysis.R} (Tier A) errors on
\texttt{read.vcfR("merged\_cohort.vcf")} & Working directory is not
\texttt{WORKDIR/variants}, or the file wasn't regenerated after Step 6's
re-merge & \texttt{cd} into \texttt{WORKDIR/variants} first; confirm the
intended cohort VCF (original vs.~haploid-corrected) actually exists
there under the name the script expects \\
Script 7's \texttt{VARIANT\_REPORT.md} shows ``Not available'' for a
section & The corresponding earlier script was not run, or was run in a
different working directory than the one passed to Script 7 & Re-run the
missing step with the same \texttt{WORKDIR} argument, then re-run Script
7 \\
Script 7's Section 8 is blank & Expected --- intentionally left for the
analyst to complete by hand & Fill it in based on the actual
mapping/coverage flags (Section 2), Ti/Tv (Section 5), and clustering
results (Section 6) --- not the auto-generated prompt text \\
\end{longtable}
}

\subsection{Completion checklist (portfolio
artifact)}\label{completion-checklist-portfolio-artifact}

The curriculum's Module 4 assessment asks for: \emph{``A filtered VCF +
a QC summary + an explicit statement of the call set's limitations
(coverage gaps, reference bias, filter thresholds used).''} Confirm each
is present before considering the module complete:

\begin{itemize}
\tightlist
\item[$\square$]
  \textbf{Filtered VCF} ---
  \texttt{\textless{}SRR\textgreater{}\_filtered\_v2.vcf} (Script 3) or,
  better, the ploidy-corrected \texttt{merged\_cohort\_haploid.vcf}
  (Step 6) is retained alongside the report, not just summarized.
\item[$\square$]
  \textbf{QC summary} --- Section 5 of \texttt{VARIANT\_REPORT.md}
  reflects a real Script 5 run against your actual merged cohort VCF,
  including the Ti/Tv value and per-sample heterozygosity, not a
  placeholder.
\item[$\square$]
  \textbf{Coverage gaps / reference bias} --- Section 2 reflects a real
  Script 2 run (per-sample mapping rate and \% zero-coverage), not a
  placeholder.
\item[$\square$]
  \textbf{Filter thresholds used} --- Section 3 states the exact
  per-sample thresholds Script 3 actually computed for this cohort, not
  the naive fixed defaults alone.
\item[$\square$]
  \textbf{Ploidy status} --- Section 4 states plainly whether the
  reported calls are diploid-default or ploidy-corrected (Script 4), and
  Section 5's heterozygosity numbers are consistent with that answer.
\item[$\square$]
  \textbf{Interpretation} --- Section 8 is filled in by hand,
  referencing the actual mapping/coverage, filter, ploidy, and Ti/Tv
  findings above --- not left as the auto-generated prompt text.
\end{itemize}

\end{document}
